% Copyright statement
\pagenumbering{Roman}

\pdfbookmark[0]{Copyright Statement}{copyright}
\includepdf[pages=-]{Appendix/copyright.pdf}

% 

\pdfbookmark[0]{Front Matter}{front}
\cftaddtitleline{toc}{chapter}{Front Matter}{}

%% =============================================================== PhD Titlepage

\begin{titlepage}
\pdfbookmark[1]{Title Page}{title}
\thispagestyle{empty}
\pagenumbering{roman}

\begin{center}
\setstretchtitle

\vspace*{1.5cm}

\begin{figure}[h!]
\centering
\includegraphics[width=12cm]{Preamble/usyd_logo} % Logo
\end{figure}

\vspace{2cm}
 
\Large{
	{\sffamily
		HIGH FIDELITY BIOELECTRIC MODELLING\\
		OF THE IMPLANTED COCHLEA}
	}

	% BOUNDARY~CONDITIONS, VASCULAR~STRUCTURE,\\
	% AND TIME-DEPENDENT EFFECTS:\\
	% GUIDELINES FOR IMPROVING VOLUME CONDUCTION MODELS
	% OF THE IMPLANTED~COCHLEA
		
	% A CRITCAL REVIEW OF || ON THE IMPORTANCE OF\\
	% BOUNDARY~CONDITIONS, VASCULAR~STRUCTURE, AND TIME-DEPENDENT EFFECTS\\
	% IN VOLUME CONDUCTION MODELS\\
	% OF THE IMPLANTED~COCHLEA}
	
	% RESOLVING THE \textit{AD POPULUM} FALLACIES IN VOLUME~CONDUCTION MODELS
	% OF THE IMPLANTED~COCHLEA
	
	% THE EFFECT OF VASCULAR~STRUCTURE IN AN ELECTRO-ANATOMIC FINITE~ELEMENT~MODEL
	% OF VOLUME~CONDUCTION IN THE COCHLEA
	
% \vspace{\stretch{3}}
% \Large{
% 	{\sffamily Volume 2}
% }
	

\vspace{\stretch{3}}

\large {
	Paul Chun Hymn Wong
}

\vspace{1ex}

\normalsize{
	\submissiondate
}

\vspace{\stretch{3}}

\footnotesize{
	A thesis in conjunction with Cochlear Limited\\
	Submitted in fulfilment of the requirements for the degree of
	Doctor~of~Philosophy\\
	School~of~Aerospace, Mechanical and Mechatronic Engineering\\
	Faculty of Engineering and Information Technologies\\ 
	The~University~of~Sydney
}

\end{center}

\end{titlepage}

%% ================================================================= Declaration

\setstretchnormal

\chapter*{Declaration}
\addcontentsline{toc}{section}{Declaration}
\setcounter{page}{2}

I declare that the work in this thesis has not previously been submitted for a 
degree nor has it been submitted as part of requirements for a degree except as 
fully acknowledged within the text.

I also declare that the thesis has been written by me. Any help that I have 
received in my research work and the preparation of the thesis itself has been 
acknowledged. In addition, I declare that all information sources and literature 
used are indicated in the thesis.

% Signature of Student
\vspace{3cm}

Paul Chun Hymn Wong

\submissiondate

%% ==================================================================== Abstract

\chapter*{Abstract}
\addcontentsline{toc}{section}{Abstract}

% Abstract header/footer
	\pagestyle{fancy}
	\lhead{{\sffamily \MakeUppercase{Abstract}}}
	\chead{}
	\rhead{}
	\lfoot{}
	\cfoot{{\sffamily \thepage}}
	\rfoot{}

% From Priya Narasimhan, ``How to write a good (no, great) PhD dissertation''
% Carnegie Mellon University

% Paragraph 1---What is the problem? 3--4 sentences about what the problem is in
% simple words.
% 
% Paragraph 2---Why is the problem hard? Review some key literature and identify
% knowledge gap, then obstacles to solving it.
% 
% Paragraph 3---What is your approach to solving this problem? ``Sit up and take
% notice'' claims that your thesis plans to demonstrate.
% 
% Paragraph 4---What is the consequence of your approach? Describe the impact of
% the research and what it enables.

Cochlear implants have been highly successful at restoring sound perception in
individuals with sensorineural hearing loss. Since their inception, improvements
in sound perception have largely been driven by advances in signal processing.
Despite the merits of newer and more advanced processing schemes however, the
level of improvement has plateaued for almost a decade, and performance in some
everyday scenarios still leaves much to be desired. This suggests that there is
an information bottleneck at the electrode-tissue interface, which is ultimately
responsible for enacting the biophysical changes that govern neuronal
recruitment during electrical stimulation.

Developing a complete understanding of the electrode-tissue interface is
challenging because the cochlea is small, consists of many intricate structures,
and is difficult to access. Modelling techniques, in conjunction with \invivo{}
measurements, have provided useful insights into the electrical response of
cochlear tissues, and recent advances in computational power have enabled
spatially accurate volume conduction models for estimating the induced electric
field. The predicted field can be fed into neural excitation models to predict
the neuronal response. However, the models depend on a number of assumptions,
some of which have yet to be tested, and despite their sophistication, many
models are unable to predict patient outcomes consistently.

The present thesis aims to improve the reliability of volume conduction models
by creating new, high fidelity reconstructions of the inner ear and critically
assessing some of the hitherto untested modelling assumptions. Firstly, the lack
of consensus around boundary conditions for monopolar simulations is addressed.
The evidence suggests that the unmodelled return path should be accounted for,
perhaps by applying a voltage offset at a boundary surface some distance from
the cochlea itself. Secondly, the role of the cochlear vasculature, which has
been ignored despite early calls for investigation, is examined. The models show
that the large modiolar vessels, particularly the vein of the scala tympani, can
have a strong local effect on current flow and neural excitation near the
stimulating electrode. Thirdly, a time-dependent simulation is performed in
response to questions about the validity of the oft-cited quasi-static
assumption. It appears that the assumption is not valid due to the relatively
high permittivity of neural tissue.

By demonstrating the impact of these factors, it is hoped that the
trustworthiness of all bioelectric models of the cochlea is improved, either by
further validating the claims of existing models, or by prompting improvements
in future studies. Developing our understanding of the physics underlying
cochlear implant systems will enable more robust and efficacious implant designs
in the future, ultimately paving the way for improved sound perception and
quality of life in individuals with sensorineural hearing loss.

%% ================================================================== Dedication

\chapter*{\hphantom{Dedication}}
\thispagestyle{empty}
\pdfbookmark[1]{Dedication}{dedication}

\begin{figure}[!t]
\centering
\vspace{4cm}

	\begin{minipage}[t]{0.77\textwidth}
	\centering
	\setstretchnormal
		I dedicate this thesis to my good friend Mitchell Seow.
		
		I can't believe it has almost been a decade since your life was cut short so
		unexpectedly. I hope that by earning this degree, I am in some small way
		making up for the opportunities that you missed out on.
		
		See you on the other side.
	\end{minipage}
	
\end{figure}

%% ===================================================================== Preface

\chapter*{Preface}
\addcontentsline{toc}{section}{Preface}

% Preface header/footer
	\pagestyle{fancy}
	\lhead{{\sffamily \MakeUppercase{Preface}}}
	\chead{}
	\rhead{}
	\lfoot{}
	\cfoot{{\sffamily \thepage}}
	\rfoot{}

As someone who normally does a lot of introspection, I was surprised to feel a
little apprehensive upon looking back on the last four years of my life. What
had I achieved? Where might I be if I had followed a different path? Was it all
worth it?

I remember starting off feeling a little uncertain about the decision to pursue
the PhD. Despite being at the same institution, everything felt a little foreign
(spending my first week at an overseas conference may have contributed to this).
The friends I had completed my undergraduate studies with were no longer on
campus, and those who had stayed behind were the quiet, unassuming types who
mostly kept to themselves. My initial bunker buddy, Scott Townsend, bucked this
trend but possessed a tenacity that I did not have, adding to the doubts in my
head. The person who helped me get through that first semester though was
undoubtedly the aptly named Gabby Chan. Her frequent visits and persistent
noise-making were a welcome distraction from the punishing initiation that is
the literature review, so I'm glad that she was still around during that time.

Four months in, Scott and I were moved downstairs to the cubicle farm in S203,
otherwise known as the domain of the tissue engineers. I felt even more out of
place there, but appreciated the open space, the air conditioning, and
especially the abundance of natural light. By then I had started to realise how
much there was to do; Howard Lau reinforced this when he pointed out that the
topics Phil and I had chosen were effectively three projects in one. Given the
incredibly niche scope of my original topic, the technical difficulties of
obtaining a suitable image stack to actually begin the modelling work, and my
own lack of relevant background, the limited early progress was disheartening.
However, my search for data lead to me to cross paths with Ian Curthoys and
Christopher Wong, who were happy to collaborate on the project and kind enough
to save me the trouble of obtaining my own ethics approval.

A few months later I moved offices once again, giving up the natural light to
join Phil and Andrian in Link 243. Soon enough, that office became a space we
could call our own. Lights and air conditioning were installed, new furniture
was procured, and over time a plethora of customisations were added.
Productivity increased as we bounced ideas off each other and became more aware
of the intricacies of our topics. A high resolution image stack that seemed
ideal for my model was finally made available, though on the flipside, I was
faced with the dauntingly tedious task of semi-manually segmenting those scans
for reconstruction. In conjunction with a quarter-life crisis, 2012 was a rather
dark year. Fortunately, it was capped off by a delightful trip to Langkawi,
Malaysia, where I presented my research work for the first time and made a heap
of great memories with some new friends from Europe.

As it turned out, 2012 was not the end of the world, and by early 2013 it became
clear that I needed more balance in my life. My internal supervisor gave me the
chance to gain some official lecturing experience, which I enjoyed (though part
of me regrets spending too much time on teaching over the course of the PhD). I
started dabbling in new experiences, most notably the plunge into digital
photography. The cost of entry was particularly painful for my wallet, but has
since enriched my life many times more than the initial outlay. We saw the Dalai
Lama and witnessed his irrepressible mirth and optimism. I also got to travel
more, presenting my research at CIAP2013 by Lake Tahoe (bookended with visits to
San Francisco and New York) and Medical Bionics 2013 on Phillip Island (followed
by a few days in Melbourne). My sister commented around then that I had become
more introverted since starting the PhD. In hindsight, this was probably true,
but I was also still recovering from the events of the previous year.

2014 was a year of change. I began socialising more and started dating again. I
established a collaboration with Shefin George at the Bionics Institute and got
to supervise two undergraduate honours theses, which expanded the depth of my
thesis. My sister got married, and I shed tears in front of hundreds of wedding
guests. My external supervisor spent most of the year abroad, and his expertise
was often sorely missed as I prepared my first manuscript for publication.
Towards the end of the year, I won an award at the Student Research Conference,
and started seeing the end point\ldots or so I thought until feedback on my
paper submission prompted me to work on additional extensions to my model.

From the start of 2015 onwards, things were a blur. Between the model/manuscript
revisions, thesis writing, and another visit to CIAP (and Seattle/Vancouver),
there was barely any time for socialising or sleeping.%
\footnote{\href{https://allthingsstick.files.wordpress.com/2012/02/college2.jpg?w=690}{This
meme} summarises the situation pretty well.} Now that I've reached the end of
this journey though, I'm glad to have compiled a thesis that I am proud to
submit. I've learnt a lot over this time, not only in terms of technical
knowledge but also a variety of soft skills that I didn't expect to pick up when
I started. Perhaps most importantly, I have discovered that I really enjoy
learning new things. And exploring the world.

To all newcomers to the PhD world, chances are you'll end up making similar
mistakes as I did (and consequently learning the same lessons), so for what it's
worth, I'd like to offer a few words of advice based on my own experience:

\begin{description}
  	\item[\textsf{Multithread your work}] The first thing you should do is
  	figure out what needs to be done using, for instance, a Gantt chart. You'll
  	inevitably find that some tasks can be done independently of others. Use
  	that to optimise your workflow because each task will almost certainly take
  	longer to achieve than you have planned (such is the nature of research).
  	When you hit a roadblock on one task, having an alternative to move on to
  	will help you to maintain your momentum and give your mind time to come up
  	with a solution.
  	\item[\textsf{Use a productivity system}] You may have gotten away without
  	using one of these before, but following one will help you to stay on top of
  	your tasks. There are many different systems out there, so try them out early
  	on to discover which works best for you. The better systems track your to-dos
  	as well as what you've already done to remind you of your accomplishments.
  	For me, \href{https://habitrpg.com/}{HabitRPG} was best for planning and tracking,
  	while an ambient noise generator (I settled on
  	\href{http://www.noisli.com/}{Noisli} for a mix of cafe sounds and brown
  	noise), earbuds, and a good Pomodoro timer were a terrific combination for
  	really knuckling down. (I later wrote my own Python script to combine these
  	latter components into one convenient package---I called it
  	``Paulmodoro''.) Don't be afraid to rest when you need to though---working
  	while fatigued is usually counterproductive.
	\item[\textsf{Learn to use \LaTeX}] Microsoft Word is adequate, but
	unprofessional and prone to errors. \LaTeX{} was designed precisely for writing
	up long documents and makes it significantly easier to keep your dissertation
	neat and organised, even if your source files are not. Combining \LaTeX{} with
	AutoHotkey also makes inserting references ridiculously simple. The benefits
	manifest most clearly in the later stages of writing up when your image and
	citation counts are in the hundreds, so it's definitely worth investing some
	time early on to set up your editor and learn the ropes.
  	% TODO: Put up a template somewhere?
	\item[\textsf{Find good collaborators}] Your PhD is a pimple on a sliver
	in the vastness of all human knowledge.%
	\footnote{As depicted brilliantly and succinctly by
	\href{http://matt.might.net/}{Matt Might} in
	\href{http://matt.might.net/articles/phd-school-in-pictures/}{The Illustrated
	Guide to a Ph.D.}.}
	While you may be tempted to learn and do everything yourself, you'll soon
	realise that it's much more efficient to collaborate with someone who already
	has expertise in that area, especially if it's not a major component of your
	project. They'll make your life a lot easier, whether it's by offering skills
	that would take you months or years to master yourself, or letting you
	piggyback on studies that already have ethics approval. One way to meet
	potential collaborators is to\ldots
	\item[\textsf{Go to conferences}] Travelling (either domestically or
	internationally) is a wonderful and enriching experience. Apply for financial
	assistance and present your work to other researchers because you'll receive
	valuable feedback while also learning more about your field. It's not all about
	the technical matters though---meeting people from all over the world is a
	great way to make new friends and reminds you that you're not alone in the
	quest for knowledge.
	\item[\textsf{Remember the bigger picture}] Your PhD may feel omnipresent and
	all-consuming at times, but you do (or should) have a life outside of the PhD
	as well. Don't neglect those things, and try to keep everything in balance as
	much as possible. Your relationships with those who matter most in your life
	are especially important. There will inevitably be times when you will stumble
	and fall, and those people are the ones who will help you pick yourself up
	again.
\end{description}

Finally, a closing thought:

\begin{verse}
	\textit{
		``Happiness does not come from doing easy work \hphantom{here is some
		more filler text} but from the afterglow of satisfaction that comes after the
		achievement of a difficult task that demanded our best.'' }

	\raggedleft{
		--- Theodore Isaac Rubin
	}
\end{verse}

Embrace the challenge and make sure you submit something that you are proud of
because it will make this accomplishment that much more fulfilling when you look
back on it in the years to come.

%% ============================================================ Acknowledgements

\chapter*{Acknowledgements}
\addcontentsline{toc}{section}{Acknowledgements}

% Acknowledgements header/footer
	\pagestyle{fancy}
	\lhead{{\sffamily \MakeUppercase{Acknowledgements}}}
	\chead{}
	\rhead{}
	\lfoot{}
	\cfoot{{\sffamily \thepage}}
	\rfoot{}

First and foremost, I would like to thank my supervisors: Qing Li and Paul
Carter. This project would not have been possible without your combined and
generous guidance. To Qing, your commitment to my personal development and
unwavering support over the past four years have demonstrated to me the true
value of a good mentor, and I will aspire to emulate these qualities in my
future endeavours. To Paul, your unfathomable technical knowledge and
first-principles understanding of all things related to cochlear implant design
have inspired me since we first met as part of the Advanced Engineering program.
I appreciate the opportunities that you have afforded me and hope that our
friendship continues to grow in the years to come.

To Ian Curthoys, Christopher Wong, and Allan Jones, thank you for being such
willing collaborators and for sharing various microCT image stacks with me
over the course of this project. I have always felt humbled by your kindness
and expertise, and your work has opened my eyes to research outside my field.

Thank you to Peter Santi and Jonathon Kirk. Although I have not had the chance
to meet either of you in person, I am very grateful for the sTSLIM dataset
that was made available for this project and hope that the model made using
those images helps to showcase the importance of your own research.

To Nick Pawsey and Joe Giorgio, thank you for providing access to Amira, for
generating CAD models of the electrode arrays, and for your ongoing support and
discussion regarding the various modelling issues. I hope that my work here adds
some value to your ongoing efforts at Cochlear.

I would also like to thank Rob Shepherd and Shefin George from the Bionics
Institute in Melbourne, Australia. Rob, thank you for showing such interest in
my work at CIAP2013, and for encouraging me to submit an abstract for Medical
Bionics 2013. Shefin, thanks for being such an awesome collaborator and for
sharing your experimental data with me. I know you're going to go far and wish
you all the best in your career.

Thank you to Payal Mukherjee and Bill Gibson for allowing me to attend a
cochlear implant surgery. As a biomedical engineer with a fascination for
human anatomy, I have always been excited to see a procedure in person.
Witnessing the clinical context first-hand has allowed me to gain a more
holistic understanding of all the effort required to complete the journey of
a medical implant, and for that I am very grateful.

A special shout out to my undergraduate honours students, Li Cen and Shilpitha
Inguva. Your friendship and assistance during 2014 was invaluable and I hope
that we can all meet up again at some point to enjoy a game of mini-pool.

To my office buddies Scott Townsend, Phillip Tran, and Andrian Sue, I have
appreciated your mateship and encouragement throughout my postgraduate life
more than you know. To Phil and Andrian in particular, I would not have gotten
through this whole ordeal without your ongoing support. You've both taught me
a lot during our time together in terms of both technical skills and more
general life advice. Regardless of where our future paths take us, I know we
will all look back on our many adventures during this time fondly.

I also could not have gotten through this without the encouragement of my
closest friends (you know who you are). As for a few special mentions, I'd like
to acknowledge Gabby Chan for her company, especially during the early months of
my PhD life. Thanks to Calla Klafas and Carina Blaker for their many hours of
open conversation over the past few years. Thanks also to Lulu Tan for
introducing me to Melinda Nguyen.

Lastly and most importantly, I would like to thank my family for their
unconditional love and support. I feel extremely privileged to have had the
opportunity to undertake study at such a high level. Thank you for being patient
with me, for putting up with my bad habits, and for being there through thick
and thin. I love you all dearly.

\clearpage

%% ==================================================================== Contents

% Separate formatting
\setstretchtoc

% TOCLOFT header/footer
	\pagestyle{fancy}
 	\lhead{{\sffamily \MakeUppercase\leftmark}}
	\chead{}
	\rhead{}
	\lfoot{}
	\cfoot{{\sffamily \thepage}}
	\rfoot{}
	
\setlength{\cftbeforetoctitleskip}{91pt}
\setlength{\cftbeforeloftitleskip}{91pt}
\setlength{\cftbeforelottitleskip}{93pt}

\phantomsection
\pdfbookmark[1]{Table of Contents}{contents}
\tableofcontents
\pagebreak

\phantomsection
\addcontentsline{toc}{section}{List of Figures}
\listoffigures
\pagebreak

\phantomsection
\addcontentsline{toc}{section}{List of Tables}
\listoftables

% Revert formatting
\setstretchnormal

%% ======================================================== List of Publications

\chapter*{List of Publications}
\addcontentsline{toc}{section}{List of Publications}

% LOP header/footer
	\pagestyle{fancy}
	\lhead{{\sffamily \MakeUppercase{List of Publications}}}
	\chead{}
	\rhead{}
	\lfoot{}
	\cfoot{{\sffamily \thepage}}
	\rfoot{}

The following is a complete listing of published works arising from this thesis.
Copies of the papers are provided in Appendix~\ref{appendix:digital_files}.

% Get formatting from tmp/Thesis.bbl

\section*{Peer-reviewed Journal Papers}
\begin{itemize}
	\item P.~Tran, A.~Sue, \textbf{P.~Wong}, Q.~Li, and P.~Carter, ``Development of
	HEATHER for cochlear implant stimulation using a new modeling workflow,''
    \emph{Biomedical Engineering, IEEE Transactions on}, vol.~62, no.~2, pp.
    728--735, February 2015.
	\item \textbf{P.~Wong}, S.~George, P.~Tran, A.~Sue, P.~Carter, and Q.~Li,
	``Development and Validation of a High-Fidelity Finite-Element Model of
	Monopolar Stimulation in the Implanted Guinea Pig Cochlea,'' \emph{Biomedical
	Engineering, IEEE Transactions on}, vol.~63, no.~1, pp. 188-198, January 2016.
\end{itemize}

\section*{Peer-reviewed Full Length Conference Papers}
\begin{itemize}
	\item \textbf{P.~Wong}, Q.~Li, and P.~Carter, ``Incorporating vascular
	structure into electric volume conduction models of the cochlea,'' in
	\emph{Biomedical Engineering and Sciences (IECBES), 2012 IEEE EMBS Conference
	on}.\hskip 1em plus 0.5em minus 0.4em\relax IEEE, December 2012, pp. 694--699.
	\item P.~Tran, \textbf{P.~Wong}, A.~Sue, Q.~Li, and P.~Carter, ``Influence of
	blood vessel conductivity in cochlear implant stimulation using an finite
	element head model,'' \emph{Engineering in Medicine and Biology Society (EMBC),
	35th Annual International Conference of the IEEE}.\hskip 1em plus 0.5em minus
	0.4em\relax IEEE, July 2013.
	\item A.~Sue, P.~Tran, \textbf{P.~Wong}, Q.~Li, and P.~Carter, ``Time-domain
	finite element models of electrochemistry in intracochlear electrodes,''
	\emph{Engineering in Medicine and Biology Society (EMBC), 35th Annual
	International Conference of the IEEE}.\hskip 1em plus 0.5em minus 0.4em\relax
	IEEE, July 2013.
	\item C.~Inguva, \textbf{P.~Wong}, A.~Sue, A.~McEwan, and P.~Carter,
	``Frequency-Dependent Simulation of Volume Conduction in a Linear Model of
	the Implanted Cochlea,'' \emph{Neural Engineering (NER), 7th International
	IEEE EMBS Conference on}.\hskip 1em plus 0.5em minus 0.4em\relax IEEE,
	April 2015.
	\item A.~Sue, \textbf{P.~Wong}, P.~Tran, Q.~Li, and P.~Carter,
	``Modeling the Effects of Electrode Recessing on Electrochemical Safety in
	Cochlear Implant Electrodes,'' \emph{Neural Engineering (NER), 7th
	International IEEE EMBS Conference on}.\hskip 1em plus 0.5em minus 0.4em\relax
	IEEE, April 2015.
\end{itemize}

Selected as an Open Finalist at NER2015 (Inguva~\etal).

\section*{Peer-reviewed Conference Abstracts with Posters}
\begin{itemize}
	\item \textbf{P.~Wong}, P.~Tran, A.~Sue, Q.~Li, and P.~Carter, ``Sensitivity of
	cochlear volume conduction models to boundary conditions,''
	\emph{Implantable Auditory Prostheses (CIAP), 2013 Conference on}.\hskip 1em
	plus 0.5em minus 0.4em\relax ARO, July 2013.
	\item P.~Tran, \textbf{P.~Wong}, A.~Sue, Q.~Li, and P.~Carter, ``Determining
	the Boundary Conditions for Cochlear Implant Models,'' \emph{Implantable
	Auditory Prostheses (CIAP), 2013 Conference on}.\hskip 1em plus 0.5em minus
	0.4em\relax ARO, July 2013.
	\item \textbf{P.~Wong}, A.~Sue, P.~Tran, Q.~Li, and P.~Carter, ``Effect of
	petrous bone resistivities on volume conduction in the cochlea,'' \emph{Medical
	Bionics, 3rd International Conference on}.\hskip 1em plus 0.5em minus
	0.4em\relax The Bionics Institute, November 2013.
	\item A.~Sue, \textbf{P.~Wong}, P.~Tran, Q.~Li, and P.~Carter, ``Modelling the
	Electrochemistry of Platinum Intracochlear Stimulation Electrodes,''
	\emph{Medical Bionics, 3rd International Conference on}.\hskip 1em plus 0.5em
	minus 0.4em\relax The Bionics Institute, November 2013.
	\item \textbf{P.~Wong}, A.~Sue, P.~Tran, C.~Inguva, Q.~Li, and P.~Carter,
	``Time-domain simulation of volume conduction in the guinea pig cochlea,'' in
	\emph{Implantable Auditory Prostheses (CIAP), 2015 Conference on}.\hskip 1em
	plus 0.5em minus 0.4em\relax ARO, July 2015.
	\item P.~Tran, \textbf{P.~Wong}, A.~Sue, Q.~Li, and P.~Carter, ``A multi-scale
	model of cochlear implant stimulation,'' in \emph{Implantable Auditory Prostheses
	(CIAP), 2015 Conference on}.\hskip 1em plus 0.5em minus 0.4em\relax ARO, July
	2015.
	\item A.~Sue, P.~Tran, \textbf{P.~Wong}, Q.~Li, and P.~Carter, ``Investigations
	of irreversible charge transfer from cochlear implant electrodes: in vitro and
	in silico approaches,'' in \emph{Implantable Auditory Prostheses (CIAP), 2015
	Conference on}.\hskip 1em plus 0.5em minus 0.4em\relax ARO, July 2015.
\end{itemize}

Received a Monash Vision Group prize at Medical Bionics 2013 (Wong~\etal).

\section*{Other Presentations}
\begin{itemize}
	\item \textbf{P.~Wong}, ``The effect of vascular structure on volume conduction
	in the cochlea," \emph{Cochlear International R\&D Symposium}.\hskip 1em plus
	0.5em minus 0.4em\relax Cochlear Limited, June 2012.
	\item \textbf{P.~Wong}, ``Modelling electrical stimulation in the implanted
	cochlea," \emph{Cochlear internal modelling workshop}.\hskip 1em plus 0.5em
	minus 0.4em\relax Cochlear Limited, May 2014.
  	\item \textbf{P.~Wong}, ``Bioelectric modelling of the implanted guinea pig
	cochlea," \emph{6th Annual Student Research Conference}.\hskip 1em plus 0.5em
	minus 0.4em\relax The University of Sydney, November 2014.
\end{itemize}

Awarded ``Best Presentation from the School of AMME'' at the 2014 Student
Research Conference. Smaller meetings with the computational modelling team at
Cochlear Limited were also held on numerous occasions.

%% =============================================================== Abbreviations

% \chapter*{Abbreviations}
% 
% \section*{Common Terms}
% 
% \begin{description}
% 
% \item[BEM] Boundary Element Method
% \item[CAP] Compound Action Potential
% \item[CHL] Conductive Hearing Loss
% \item[CI] Cochlear Implant
% \item[CT] Computed Tomography
% \item[eCAP] Electrically-evoked Compound Action Potential
% \item[FDM] Finite Difference Method
% \item[FEM] Finite Element Method (or Model)
% \item[HR] High Resolution
% \item[LSFM] (Laser) Light-Sheet-based Fluorescence Microscopy
% \item[N VIII] 8th Cranial Nerve (i.e. the auditory nerve)
% \item[OC] Organ of Corti
% \item[OPFOS] Orthogonal-Plane Fluorescence Optical Sectioning
% \item[OW] Oval Window
% \item[ROI] Region Of Interest
% \item[RW] Round Window
% \item[SHL] Sensorineural Hearing Loss
% \item[TSLIM] Thin-Sheet Laser Imaging Microscopy
% 
% \end{description}
% 
% 
% \section*{Blood Vessels}
% 
% \begin{description}
% 
% \item[AVA] Anterior Vestibular Artery
% \item[SMA] Spiral Modiolar Artery
% \item[VCA] Vestibulo-Cochlear Artery
% 
% \item[CMV] Common Modiolar Vein
% \item[VCAQ] Vein of the Cochlear Aqueduct
% \item[VCV] Vestibulo-Cochlear Vein
% \item[VST] Vein of the Scala Tympani
% \item[VSV] Vein of the Scala Vestibuli
% 
% \end{description}

%% =================================================================== Quotation

\chapter*{\hphantom{Quotation}}
\thispagestyle{empty}

\vfill

\begin{verse}
	
	\textit{
		``\ldots it is necessary to investigate the electrical geometry of the
		cochlea, because only with a knowledge of the voltage and resistance
		patterns is it possible to form an over-all picture of the electrical
		situation in the cochlea.''
	}
	
	\vspace{4mm}
	
	\raggedleft{
		--- Georg von \bekesy, 1951
	}
	
	% Georg von B�k�sy received the Nobel Prize in Physiology or Medicine 1961 "for
	% his discoveries of the physical mechanism of stimulation within the cochlea".

\end{verse}

\vspace{1cm}

% Start page numbering at 1 on next page
\pagenumbering{arabic}
\setcounter{page}{0}